% !TEX program = lualatex
\documentclass[aspectratio=43,professionalfonts,handout]{beamer}
\usepackage{beamer_template}

\newcommand{\LectureNum}{23}
\newcommand{\LectureTitle}{加法定理}
\newcommand{\TextPages}{p.\,158-159}
\newcommand{\NextTopic}{2倍角・半角の公式}
\newcommand{\NextPages}{p.\,160-162}
\newcommand{\CourseName}{基礎数学B}
\renewcommand{\TermName}{後期}

\begin{document}

% -----------------------------------------------
% スライド1: 表紙＋目標＋用語・定義
% -----------------------------------------------
\begin{frame}{\CourseName\quad 第\LectureNum 回「\LectureTitle 」}
  {\small \TermName\quad \TeacherName\quad ／\quad 教科書 \TextPages\quad
  \textbf{目標}：$\sin$・$\cos$・$\tan$ の加法定理を理解し証明できる}

  \begin{block}{用語・定義}
  \begin{description}[labelwidth=5.5em]
    \item[\kw{$\sin(\alpha \pm \beta) = \sin\alpha\cos\beta \pm \cos\alpha\sin\beta$}]
      $\sin(\alpha+\beta) = \sin\alpha\cos\beta + \cos\alpha\sin\beta$，$\sin(\alpha-\beta) = \sin\alpha\cos\beta - \cos\alpha\sin\beta$
    \item[\kw{$\cos(\alpha \pm \beta) = \cos\alpha\cos\beta \mp \sin\alpha\sin\beta$}]
      $\cos(\alpha+\beta) = \cos\alpha\cos\beta - \sin\alpha\sin\beta$，$\cos(\alpha-\beta) = \cos\alpha\cos\beta + \sin\alpha\sin\beta$
    \item[\kw{$\tan(\alpha \pm \beta) = (\tan\alpha \pm \tan\beta)/(1 \mp \tan\alpha\tan\beta)$}]
      $\tan(\alpha+\beta) = \dfrac{\tan\alpha+\tan\beta}{1-\tan\alpha\tan\beta}$，$\tan(\alpha-\beta) = \dfrac{\tan\alpha-\tan\beta}{1+\tan\alpha\tan\beta}$
  \end{description}
  \end{block}
\end{frame}

% -----------------------------------------------
% スライド2: 公式・性質
% -----------------------------------------------
\begin{frame}{公式・性質}
  \begin{block}{}
  \begin{itemize}
    \item 加法定理6公式（$\sin$・$\cos$・$\tan$ の $+$ と $-$）
    \item 代表的な応用：$\sin 75°$、$\cos 15°$ などの計算
  \end{itemize}
  \end{block}
\end{frame}

% -----------------------------------------------
% スライド3: 例1→問1
% -----------------------------------------------
\begin{frame}{例1→問1}
  \begin{exampleblock}{例1) p.174（$\sin 75°$、$\cos 15°$ を求める）}
    （板書で解説）
  \end{exampleblock}

  \vfill

  \begin{block}{問1) p.174\quad 自分で解いてみよう}
    （ノートに解くこと）
  \end{block}
\end{frame}

% -----------------------------------------------
% スライド4: 例2→問2
% -----------------------------------------------
\begin{frame}{例2→問2}
  \begin{exampleblock}{例題1) p.175（$\sin(\alpha+\beta)$ を求める）}
    （板書で解説）
  \end{exampleblock}

  \vfill

  \begin{block}{問2) p.175\quad 自分で解いてみよう}
    （ノートに解くこと）
  \end{block}
\end{frame}

% -----------------------------------------------
% まとめ
% -----------------------------------------------
\begin{frame}{まとめ}
  \begin{itemize}
    \item $\sin(\alpha\pm\beta)$，$\cos(\alpha\pm\beta)$，$\tan(\alpha\pm\beta)$ の3セットを確実に覚える
    \item $15° = 45° - 30°$，$75° = 45° + 30°$ など特殊角への応用
    \item 加法定理から2倍角・半角公式が導出できる
  \end{itemize}

  \vfill
  {\small 基本問題：318, 319, 320}
\end{frame}

\end{document}
